\documentclass[11pt]{article}
\usepackage{amsmath,amssymb,amsthm}
\newtheorem{lemma}{Lemma}
\newtheorem{theorem}{Theorem}
\newcommand{\E}{\mathbb E}
\newcommand{\Tr}{\operatorname{Tr}}
\title{Same-measure local exponential moments on odd Wilson tori}\author{}\date{}
\begin{document}\maketitle
This note removes the dyadic compatibility restriction in the local-moment application of
the local moments fixed depth calculation, without changing the torus or the measure. It uses the exact positive Wilson
transfer operator and an elementary trace inequality. The partition-power bounds remain
those evaluated from manuscript Lemma 10.1. The source is still XP3.

\begin{lemma}[Positive trace Holder inequality]
For positive trace-class operators $A,B$ and an integer $n\ge2$,
\[
 \Tr(AB^{n-1})\le(\Tr A^n)^{1/n}(\Tr B^n)^{(n-1)/n}.
\]
\end{lemma}
\begin{proof}
After rescaling, suppose both traces on the right equal one. In an eigenbasis of $B$,
write $a_i=\langle i,Ai\rangle$, $b_i$ for its eigenvalues. Scalar Young's inequality gives
$a_ib_i^{n-1}\le a_i^n/n+(n-1)b_i^n/n$.
Spectral Jensen gives $a_i^n\le\langle i,A^ni\rangle$. Sum and then undo the rescaling.
Finite-rank approximation or monotone spectral sums give the trace-class assertion.
No parity condition on $n$ occurs.
\end{proof}

\paragraph{The exact positive Wilson transfer operator.}
Choose any coordinate as time and let $\mathcal X=G^{E_{\rm spatial}}$ be a spatial-slice
link configuration. Work in $L^2(\mathcal X,d\lambda)$ with normalized product Haar.
Let $P$ be orthogonal averaging over site gauge transformations on this slice.
For nonnegative couplings $\lambda_e$, define the product convolution kernel
\[
 Q_{\boldsymbol\lambda}(X,Y)=\prod_{e\in E_{\rm spatial}}
 k_{\lambda_e}(X_eY_e^{-1}),\qquad
 k_\lambda(U)=e^{-\lambda(N-\Re\operatorname{tr}U)}.
\]
This operator is positive. Indeed
$e^{\lambda\Re\chi}=\sum_{r,s\ge0}(\lambda/2)^{r+s}\chi^r\overline\chi^s/(r!s!)$;
each tensor-product character decomposes with nonnegative integer multiplicities.
Thus every central Fourier coefficient of $k_\lambda$ is nonnegative, and convolution
has nonnegative eigenvalues. Uniform convergence is immediate from $|\chi|\le N$.
Smoothness on the compact group makes the convolution trace class; alternatively its
positive spectral sum is $k_\lambda(I)=1$. The product kernel also has trace one.

Centrality makes $Q_{\boldsymbol\lambda}$ commute with the simultaneous site gauge
projection $P$, even when couplings vary from spatial link to spatial link.
Let $M$ be multiplication by $e^{-\beta S_{\rm spatial}(X)/2}$ and set
\[
 T_{\boldsymbol\lambda}=M Q_{\boldsymbol\lambda}P M.
\]
Both $M$ and $Q$ commute with $P$, so this is positive and trace class.
For uniform $\lambda_e=\beta$ write $T$.
Expanding the trace of $T^n$ in its kernel and the gauge averages gives exactly the original
Wilson torus partition function $Z_{n,\beta}$: the slice variables are spatial links,
the gauge variables between slices are the original temporal links, $Q$ supplies all
space-time plaquettes, and the two adjacent factors of $M$ supply each spatial plaquette
once. Every Haar factor is normalized, and $k_\lambda$ includes its original constant
$e^{-N\lambda}$. No residual gauge volume or Polyakov-loop factor has been omitted;
the full periodic product of gauge averages retains those temporal integrations.

\paragraph{Dissemination in a transverse slice direction.}
Suppose an insertion is $D(X)=e^{t\beta\sum_{p\in A}s_p(X)}$, with the plaquettes in
one spatial slice and $D$ gauge invariant. Its original numerator is $\Tr(DT^n)$.
Put $A=T^{1/2}DT^{1/2}$ and $B=T$ in the lemma. Cyclicity gives
\[
 \Tr(DT^n)\le[\Tr(D^{1/2}TD^{1/2})^n]^{1/n}[\Tr T^n]^{(n-1)/n}.
\]
Here the first trace is the exact partition function with the same insertion in EVERY
time slice: $D$ commutes with $M$ and $P$, so it changes the spatial potential by the
specified amount. The notation means $\Tr((D^{1/2}TD^{1/2})^n)$.
Dividing by the original partition function disseminates the slice insertion with exponent
$1/n$, preserving all original normalizers.

\paragraph{Dissemination in an in-plane slab direction.}
If the marked plaquettes cross one temporal slab, multiplying their Boltzmann weights by
$e^{t\beta s_p}$ replaces the corresponding temporal convolution coefficients by
$(1-t)\beta$, leaving all other coefficients at $\beta$.
For $0\le t<1$, the modified operator $T_1=M Q_{\boldsymbol\lambda}PM$ is still positive.
The original numerator is exactly $\Tr(T_1T^{n-1})$.
The lemma bounds it by $(\Tr T_1^n)^{1/n}(\Tr T^n)^{(n-1)/n}$.
The first trace repeats the SAME modified slab in every time position. Spatial potentials,
Haar measures and gauge projections have not changed.

\begin{theorem}[Local moment, every integer torus size]
For the original periodic four-dimensional $SU(N)$ Wilson measure on an $n^4$ lattice,
$n\ge3$, every plaquette $p$ and $0\le t<1$ satisfy
\begin{equation}\label{main}
 \E_{n,\beta}e^{t\beta s_p}
 \le\left(\frac{Z_{n,(1-t)\beta}}{Z_{n,\beta}}\right)^{1/n^4}.
\end{equation}
There is no evenness, dyadic, or blocking-factor assumption.
\end{theorem}
\begin{proof}
Let $p$ have orientation $(1,2)$. Apply slice dissemination in direction $3$, then in
direction $4$. At each application the currently marked plaquettes all lie in a single
slice of the chosen direction. The resulting marked set contains all translates in
coordinates $3,4$, with fixed base coordinates $1,2$.
Now apply slab dissemination with direction $1$ as time. The marked set consists precisely
of the orientation $(1,2)$ plaquettes in one temporal slab whose spatial direction-2 link
has the chosen coordinate-2 base. The modified convolution operator is positive as proved
above. Its repetition marks all coordinate-1 translates.
Finally use direction $2$ as time. All marked plaquettes are in one slab and all its
spatial direction-1 links are marked; slab dissemination marks every orientation $(1,2)$
plaquette on the torus. Four applications give exponent $1/n^4$ and the ratio of this
anisotropically reduced-coupling partition function to the original one.
Since all remaining plaquette actions are nonnegative, reducing their couplings also to
$(1-t)\beta$ only increases the numerator. This proves \eqref{main}.
Every step used the original uniform measure as denominator and an exact marked numerator;
no successive change of probability or torus size was made.
\end{proof}

\paragraph{Quantitative consequence in AI normalization.}
Put $B=N\beta$ and $d_G=N^2-1$. The manuscript's partition-power bounds give
\[
 \E e^{t\beta s_p}\le
 c_u^{R_n/n^4}c_l^{-Q_n/n^4}(1-t)^{-d_GR_n/(2n^4)}B^{2d_G/n},
 \quad Q_n=3n^4+1,\ R_n=3n^4-4n^3+1.
\]
Thus the finite-size factor is $B^{2d_G/n}$, with no chessboard-cell volume and no parity
restriction. Under the ORIGINAL physical regime $n\ge c_*\log(2+B)$ it is uniformly
bounded. In particular odd $L$ and odd $n=L^k$ cause no change in measure or sequence.
For all $n\ge3$, a tail integration also gives
\[
 \|\beta s_p\|_m\le C_m\left(1+\frac{\log(2+B)}n\right).
\]
This last estimate retains the finite-size cost instead of claiming uniformity in $\beta$
at a fixed torus size.

Consequently every conclusion of the local moments fixed depth calculation holds on the original odd-torus physical sequence
whenever its explicit size/coupling condition holds: local row moments, log-free fixed-depth
conditional covariances under actual independent effective layers, and improved first-layer
mean derivatives. The row-overlap constants still grow with depth, and the stopped inherited
mean covariance remains. This result is a local moment theorem, not an all-history MT36,
separated-decay, vacuum comparison, or Yang--Mills existence proof.
\end{document}
